% Original analytic teaching diagram, 2026-11-12.
% (a) F_{X,Y}(x,y) is the probability of the lower-left quadrant;
% (b) rectangle probability from the four corner values with signs + - - +.
\definecolor{VB}{HTML}{7188A3}\definecolor{VBF}{HTML}{E5EAF0}
\definecolor{VG}{HTML}{829C8B}\definecolor{VGF}{HTML}{E8EEE9}
\definecolor{VY}{HTML}{B4A16F}\definecolor{VYF}{HTML}{F1ECDD}
\definecolor{VR}{HTML}{A7787D}\definecolor{VRF}{HTML}{F0E5E6}
\definecolor{VP}{HTML}{9588A5}\definecolor{VPF}{HTML}{ECE8F0}
\definecolor{VI}{HTML}{3B4252}\definecolor{VM}{HTML}{6B7280}\definecolor{VL}{HTML}{4C4D4F}
\begin{tikzpicture}[x=1mm,y=1mm,>=Stealth,font=\figurefont\fontsize{8}{10}\selectfont,
lab/.style={text=VM,font=\figurefont\fontsize{7.5}{9}\selectfont}]
% ---------- (a) ----------
\fill[VBF] (-2,-2) rectangle (26,22);
\draw[VB,line width=.8pt] (-2,22) -- (26,22) -- (26,-2);
\draw[->,VL,line width=.7pt] (-4,6) -- (42,6) node[right,text=VI] {$x$};
\draw[->,VL,line width=.7pt] (6,-4) -- (6,36) node[above,text=VI] {$y$};
\draw[->,VB,line width=.8pt] (2,15) -- (-5,15);
\draw[->,VB,line width=.8pt] (18,2) -- (18,-5);
\fill[VR] (26,22) circle (0.9);
\node[text=VI,above right] at (26,22) {$(x,y)$};
\node[text=VI] at (16,13) {$F_{X,Y}(x,y)$};
\node[text=VI] at (19,-12) {(a)};
% ---------- (b) ----------
\begin{scope}[shift={(56,0)}]
\fill[VBF] (14,10) rectangle (34,26);
\draw[VB,line width=.8pt] (14,10) rectangle (34,26);
\draw[->,VL,line width=.7pt] (0,0) -- (46,0) node[right,text=VI] {$x$};
\draw[->,VL,line width=.7pt] (0,0) -- (0,36) node[above,text=VI] {$y$};
\draw[VL,dashed,line width=.5pt] (14,0) -- (14,10) (34,0) -- (34,10) (0,10) -- (14,10) (0,26) -- (14,26);
\node[text=VI,below] at (14,0) {$a_1$};
\node[text=VI,below] at (34,0) {$b_1$};
\node[text=VI,left] at (0,10) {$a_2$};
\node[text=VI,left] at (0,26) {$b_2$};
\foreach \p/\s/\pos in {{(34,26)}/+/above right,{(14,10)}/+/below left,{(14,26)}/-/above left,{(34,10)}/-/below right}{
  \fill[VR] \p circle (0.9);
  \node[text=VI,\pos,font=\figurefont\fontsize{10}{12}\selectfont] at \p {$\s$};
}
\node[text=VI] at (24,18) {$\mathbb{P}(R)$};
\node[text=VI] at (23,-12) {(b)};
\end{scope}
\end{tikzpicture}
